\subsection{Dependencies in the proof}\label{sec:dependencias}
Table~\ref{tab:dependencias} gives an order in which the estimates can be verified. In particular, the
change of law that transfers the arrival bound uses an estimate already obtained under the adaptive law.
\begin{center}
\small
\begin{tabular}{@{}L{.30\linewidth}L{.63\linewidth}@{}}
\toprule Result & Inputs and subsequent use\\ \midrule
Consistency and persistence & The results from~\cite{pc,tc} in Section~\ref{sec:importados} fix the
classes, the exit rates, and the size of the row defects.\\
Arrival under \(\nu\) & Positivity (Lemma~\ref{lem:positividad-llegada}) allows us to take logarithms.
Minors, coefficients by support, and the tail (Propositions~\ref{prop:menores-llegada},
\ref{prop:soportes-llegada}, and~\ref{prop:norma-llegada}) combine with the geometry of the types to
estimate the mean under \(\nu\).\\
Rate and transfer & Lemma~\ref{lem:previa} applies the stopped identity to that mean at a single cut-off.
With the rate in hand, the stopped density transfers the bound to \(\mu\)
(Lemma~\ref{lem:llegada-integrada}).\\
Interior and boundary & In the interior we integrate the sum of the sixteen contractions
(Appendix~\ref{app:interior}). At the boundary we separate the double minors, the variation of a pair,
and the interaction of the square (Appendices~\ref{app:boundary}, \ref{app:weights},
and~\ref{app:transfer}).\\
Limit and coefficient & Identity~\eqref{eq:signed-mean} combines the two means.
Proposition~\ref{prop:cortes-uniformes} controls their errors; Lemma~\ref{lem:dos-cortes} gives the
limit. The endpoint analysis identifies the integral in Theorem~\ref{teo:principal}.\\
\bottomrule
\end{tabular}
\refstepcounter{table}\label{tab:dependencias}\par\smallskip
\textbf{Table \thetable.} Order of the results and changes of law.
\end{center}
